% Auto-generated by tools/qbe.py project-article-update.
% Do not edit this file by hand; edit the proof artifacts or article sections instead.

\section{Generated Current Cycle Status}
\label{app:generated-cycle-status}

This appendix page is generated from the latest ABEIS proof cycle.  It is a
status bridge for article maintenance, not a polished theorem proof.  Stable
mathematical claims should be copied into the main text only after the
corresponding Lean declaration, source-paper anchor, and proof-note export are
available.

\paragraph{Cycle.}
Task \texttt{QBE-AUTO-002}, cycle \texttt{1}, run
\texttt{runs/20260615-034853-QBE-AUTO-002-cycle01}.  Generated at
\texttt{2026-06-15 04:03:30}.

\paragraph{Lean status signal.}
\begin{itemize}
  \item \texttt{QuantumBlockEncoding/RobinMatrix.lean:25411:  sorry}
  \item \texttt{QuantumBlockEncoding/RobinMatrix.lean:25441:  sorry}
\end{itemize}

\paragraph{Plain-language status.}
The current GHL case study is not blocked because the paper lacks a proof sketch.  It is blocked because ABEIS must turn the circuit in the paper into an exact matrix statement and prove that the clean block entry has the coefficient claimed by the paper.  In ordinary terms, the paper says that a specific path through the circuit leaves the desired coefficient, while Lean requires the corresponding matrix entry and all unwanted branches to be proved exactly.

\paragraph{Pre-Lean verifier candidates.}
These checks are necessary-condition filters, not proofs.  A failed exact
finite check usually indicates that the current Lean target, circuit
transcript, or index map is wrong.  A passing check only means the candidate
survived a cheaper diagnostic; the final claim still requires a named Lean
theorem or an explicit contract.
\begin{itemize}
  \item \textbf{active [0,0] entry.} Fast check: exact rational matrix or path-sum evaluation.  Necessary because if the exact finite entry is not the target coefficient, the Lean equality for that entry cannot be true.  Lean still must prove: a passing check is only a counterexample filter; Lean must still prove the named entry lemma.
  \item \textbf{remaining branch vanish/cancel.} Fast check: support and path checker for the remaining backend slots.  Necessary because if an unwanted clean-branch path survives numerically, the block projection cannot match the paper target.  Lean still must prove: Lean must still prove the zero/cancellation lemma in the formal circuit semantics.
  \item \textbf{Ry boundary convention.} Fast check: symbolic 2-by-2 rotation check.  Necessary because a mismatched half-angle convention changes the boundary amplitude before any Lean tactic is relevant.  Lean still must prove: Lean must still record the convention bridge as a source-supported theorem.
  \item \textbf{sparse-access map.} Fast check: finite range/injectivity/permutation check.  Necessary because a reversible oracle cannot exist for a colliding or out-of-range finite map.  Lean still must prove: Lean must still prove the reversible extension and cleanup obligations.
  \item \textbf{coefficient oracle clean branch.} Fast check: exact finite evaluator for f(x\_i)/N\_f.  Necessary because the final block entry uses this coefficient, so a wrong clean branch invalidates the target theorem.  Lean still must prove: Lean must still prove bounds, orthogonality, and unitary completion or keep them as contracts.
\end{itemize}

\paragraph{Current proof-DAG frontier.}
\begin{itemize}
  \item Latest handoff indicates at least one assigned lower target was already compiled; upper/middle should retire stale directives before more proof search.
  \item 1. lower1 or middle prepares the next source-proof translation packet for the remaining finite normalized-block/projection bridge feeding \textasciigrave{}oneTermRobinGamma3ProductToCoefficientObligation 3 0 0\textasciigrave{}.
  \item 2. lower2 has no new Lean theorem target until that packet names exactly one declaration. If run before then, it should make no Lean edit and report the compiled conditional bridge as route memory.
  \item 3. lower3 should verify necessary conditions for the next finite normalized-block/projection bridge and keep \textasciigrave{}product\_to\_coefficient=false\textasciigrave{} until a theorem closes the root obligation.
\end{itemize}

\paragraph{Open obligation signal.}
\begin{itemize}
  \item unchanged backend-expansion raw target: Lean \textasciigrave{}oneTermRobinGamma3BoundaryBackendBranchContributionTarget\_n3.backendExpansionStatement\textasciigrave{}; class source-contract gap plus finite counterexample; status refuted; never assign unchanged
  \item source-prepared projection to slot-\textasciigrave{}2\textasciigrave{} projected product: Lean \textasciigrave{}oneTermRobinGamma3BoundarySourcePreparedProjection\_slot2\_to\_projectedBranchProduct\_n3\textasciigrave{}; class internal GHL branch/product step under \textasciigrave{}hUniform\textasciigrave{} and \textasciigrave{}hentry\textasciigrave{}; status compiled; route memory
  \item source-prepared slot-\textasciigrave{}2\textasciigrave{} product feeds fixed product map: Lean \textasciigrave{}oneTermRobinGamma3BoundarySourcePreparedProjectionSlot2Product\_feedsFixedProductMap\_n3\textasciigrave{}; class QBE-local route-retarget guard; status compiled; all product/block/final flags false
  \item fixed product-to-coefficient theorem for \textasciigrave{}(0,0)\textasciigrave{}: Lean \textasciigrave{}oneTermRobinGamma3ProductToCoefficientObligation 3 0 0\textasciigrave{}; class coefficient equality plus normalizer algebra and boundary coefficient convention; status open; future packet
  \item diagnostic raw equality route: Lean \textasciigrave{}oneTermRobinGamma3BoundaryUnitaryEntry\_eq\_backendFold\_n3\textasciigrave{}; \textasciigrave{}oneTermRobinGamma3BoundaryEvalGateMatrices\_eq\_sevenGateMatrix\_n3\textasciigrave{}; class existing diagnostic \textasciigrave{}sorry\textasciigrave{} route; status forbidden as dependency
\end{itemize}

\paragraph{Open GHL contribution obligations.}
\begin{itemize}
  \item \texttt{Uindic}: Indicator unitary; status Permutation and self-inverse helpers compile, but this is not yet the final block-encoding proof..
  \item \texttt{RyBoundary}: Boundary-controlled Ry rotations; status Active convention audit: the Ry angle convention must be source-supported before closing the boundary amplitude proof..
  \item \texttt{RobinTheorem}: One-term Robin block-encoding theorem; status Main active target: the theorem-facing Lean statement is not closed..
  \item \texttt{GammaSlices}: Gamma wavefunction slices and coefficient product; status Several finite boundary lemmas exist; the final projection/product bridge is still open..
  \item \texttt{FigRobin}: Figure 4 Robin circuit transcript; status Transcript guard compiles for visible indicator dagger and H preparation; the active seven-gate backend remains a component..
  \item \texttt{OneD}: One-dimensional Hamiltonian block encoding; status Deferred until the one-term Robin theorem closes..
  \item \texttt{MultiD}: Multidimensional block encoding; status Planned; depends on one-term and LCU abstractions..
\end{itemize}

\paragraph{Open external technical lemma obligations.}
\begin{itemize}
  \item \texttt{tl-ghl-lemma1-banded-sparse-access}: contract-only; next action Keep as explicit external contract until the exact cited construction is formalized or imported as a theorem..
  \item \texttt{tl-ghl-lemma3-sparse-amplitude}: contract-only; next action Maintain the clean-branch contract and isolate any missing unitarity proof as an external technical lemma..
  \item \texttt{tl-ghl-theorem5-piecewise-polynomial-of}: contract-only; next action Use only as a named contract in the current theorem; do not invent stronger smoothness or range assumptions..
  \item \texttt{tl-uniform-sparse-register-preparation}: contract-only; next action Keep as a theorem-facing contract unless the exact state-preparation proof is needed to close a resource theorem..
  \item \texttt{tl-ry-boundary-amplitude-convention}: obligation; next action Do not silently change the paper angle; prove the convention bridge or record the exact cited convention..
  \item \texttt{tl-qsvt-blockencoding-simulation}: paper-cited; next action Defer until the one-term Robin theorem and LCU combination are closed..
  \item \texttt{tl-clean-block-definition}: contract-only; next action Use as the final target predicate; ensure circuit entry lemmas are bridged to this semantic definition..
\end{itemize}

\paragraph{Recent typed verifier feedback.}
\begin{itemize}
  \item Leaf \texttt{conditional\_normalizer\_product\_bridge}, class \texttt{None}; next route: Use the compiled conditional bridge as route memory; next proof work must target the remaining finite normalized-block/projection bridge for oneTermRobinGamma3ProductToCoefficientObligation 3 0 0, not the retired backendExpansionStatement..
  \item Leaf \texttt{conditional\_normalizer\_product\_bridge}, class \texttt{None}; next route: Use compiled conditional bridge as route memory; next target is finite normalized-block/projection bridge for oneTermRobinGamma3ProductToCoefficientObligation 3 0 0.
  \item Leaf \texttt{conditional\_normalizer\_product\_bridge}, class \texttt{symbolic\_bridge\_gap}; next route: Prove oneTermRobinGamma3BoundarySourcePreparedSlot2Product\_normalizerEval\_n3 from the compiled source-prepared slot-2 product bridge and oneTermRobinGamma3BoundaryProductUnderContractsEval\_n3; do not revive backendExpansionStatement or SignalEntryFold..
  \item Leaf \texttt{conditional\_normalizer\_product\_bridge}, class \texttt{symbolic\_bridge\_gap}; next route: lower2 prove oneTermRobinGamma3BoundarySourcePreparedSlot2Product\_normalizerEval\_n3; do not revive backendExpansionStatement.
  \item Leaf \texttt{source\_slot2\_product\_fixed\_map\_guard}, class \texttt{stale\_leaf}; next route: prove oneTermRobinGamma3BoundarySourcePreparedSlot2Product\_normalizerEval\_n3, optionally through oneTermRobinGamma3BoundaryHWKappaUniformAllSlots\_to\_productRouteFocusedCleanColumn\_n3; do not revive backendExpansionStatement.
  \item Leaf \texttt{source\_slot2\_product\_fixed\_map\_guard}, class \texttt{stale\_leaf}; next route: Retire the source\_slot2\_product\_fixed\_map\_guard leaf as compiled route memory. Future lower work should use the new coefficient/normalizer packet for oneTermRobinGamma3ProductToCoefficientObligation 3 0 0 and must not revive backendExpansionStatement or SignalEntryFold..
\end{itemize}

\paragraph{Recent proof-attempt memory.}
\begin{itemize}
  \item \texttt{proof-attempts/QBE-AUTO-002/product-to-coefficient-normalizer-middle-synthesis-20260615-034853.md}
  \item \texttt{proof-attempts/QBE-AUTO-002/product-to-coefficient-normalizer-lower1-proof-architect-20260615-034853.md}
  \item \texttt{proof-attempts/QBE-AUTO-002/product-to-coefficient-normalizer-lower1-proof-architect-20260615-034610.md}
  \item \texttt{proof-attempts/QBE-AUTO-002/product-to-coefficient-normalizer-middle-packet-20260615-032327.md}
  \item \texttt{proof-attempts/QBE-AUTO-002/backend-expansion-route-retarget-lower-proof-architect-20260615-032040.md}
\end{itemize}

\paragraph{Article-writing rule.}
The project report may discuss this cycle as process evidence.  It must not
claim completion of the first case study \citep{guseynovHuangLiu2025} unless
the final theorem-facing block-extraction statement is closed in Lean and the
Markdown/LaTeX proof map has been synchronized.
